\documentclass[10pt,a4paper]{article}
\usepackage[utf8]{inputenc}
\usepackage[ngerman]{babel}
\usepackage[T1]{fontenc}
\title{Vereinssatzung des \glqq Die Ungezählten e.V.\grqq}
\date{01.01.2027}

\renewcommand\thesection{§\arabic{section}}

\begin{document}
\maketitle

\section{Name und Sitz}
Der Verein führt den Namen \glqq Die Ungezählten e.V.\grqq.
Sein Sitz ist in Augsburg in Bayern.
Der Verein soll in das Vereinsregister eingetragen werden.
Das Geschäftsjahr des Vereins und seiner Abteilungen ist das Kalenderjahr.

\section{Zweck, Mittelverwendung und Aufgaben des Vereins}\label{zweck}

\subsection{Zweck}
\glqq Die Ungezählten e.V.\grqq\,
verfolgt ausschließlich und unmittelbar gemeinnützige Zwecke im Sinne des Abschnitts
\glqq Steuerbegünstigte Zwecke\grqq\, der Abgabenordnung.


Der Zweck des Vereins ist die Förderung von Kunst und Kultur,
insbesondere durch die Pflege und Ausübung
von historischen und fiktionalen Rollenspielen (Live-Action-Role-Playing, LARP), der darstellenden Künste
sowie der historischen und modernen Schaustellerei. Dies umfasst sowohl den Erhalt und die Pflege des
historischen Brauchtums des Mittelalters als auch die Förderung weiterer LARP-Genres wie Fantasy,
ScienceFiction und anderer kulturell relevanter Szenarien. Der Verein fördert die Teilnahme und Inklusion aller
interessierten Personen unabhängig von Alter, Geschlecht, Herkunft sowie körperlichen und geistigen
Fähigkeiten, ohne sich auf eine spezifische Zielgruppe zu beschränken.

Der Verein erreicht seinen Zweck durch die regelmäßige Durchführung von LARP-Veranstaltungen,
die das kreative Ausdruckspotential, die kulturelle Vielfalt und die Diversität fördern. 
Zudem pflegt der Verein den Kontakt zu in- und ausländischen Vereinen und Organisationen mit ähnlicher Zielrichtung.
Dabei ist der Verein politisch und konfessionell neutral.

\subsection{Mittelverwendung}
Der Verein ist selbstlos tätig.
Er verfolgt nicht in erster Linie eigenwirtschaftliche Zwecke.
Mittel des Vereins dürfen nur für satzungsgemäße Zwecke verwendet werden.
Die Mitglieder erhalten keine Zuwendungen aus den Mitteln des Vereins.
Es darf keine Person durch Ausgaben, die dem Zweck fremd sind,
oder durch unverhältnismäßig hohe Vergütung begünstigt werden.
Die von Mitgliedern für und in den Verein eingebrachten und auch ausgearbeiteten Ideen und Projekte werden dem Verein freiwillig und selbstlos von Seiten der Mitglieder übergeben,
ebenso das Recht am geistigen Eigentum einer Idee oder eines Projekts.
Nach Austritt eines Mitgliedes bleiben dessen Ideen und Projekte,
welche es im Sinne des Vereins während seiner Mitgliedschaft erstellte,
Eigentum des Vereins und können somit von diesem weiter genutzt werden.
Der Verein unterscheidet zwischen an den Verein übergebene Ideen und Projekte und nur zur Verfügung gestellte Ideen und Projekte.
\subsection{Aufgaben}

Zu den vorrangigen Aufgaben des Vereins gehören insbesondere:
\begin{enumerate}
\item
die Ausbildung von Mitgliedern zur Teilnahme an Mittelaltermärkten und sonstigen thematisch gleichen Veranstaltungen
\item
Durchführung geeigneter Veranstaltungen für Mitglieder und Interessenten zur Förderung nach \ref{zweck}.
\end{enumerate}

\section{Mitgliedschaft}
Vereinsmitglieder können natürliche, aber auch juristische Personen werden. Jugendliche unter 18 Jahren
bedürfen der Zustimmung der/des gesetzlichen Vertreter/s.


Mitglieder des Vereins sind, soweit die Vereinsordnung keine abweichende Regelung enthält:
\begin{itemize}
\item
Ordentliche Mitglieder:
\begin{itemize}
\item Volljährige natürliche Personen
\item Jugendliche (ab der Vollendung des 16. Lebensjahres)
\item Juristische Personen
\end{itemize}
\item 
Sondermitglieder
\begin{itemize}
\item
Ehrenmitglieder
\item
Fördermitglieder
\item
Jugendliche und Kinder bis zur Vollendung des 16. Lebensjahrs
\end{itemize}
\end{itemize}


\subsection{}
Die Mitglieder sind verpflichtet, die Vereinssatzung anzuerkennen, die Zwecke des Vereins zu fördern und
zu unterstützen, die festgesetzten Mitgliedsbeiträge und Umlagen rechtzeitig zu entrichten, die Anordnungen der
erweiterten Vorstandschaft und die Beschlüsse der Mitgliederversammlung zu respektieren, den Vereinsfrieden
zu fördern und zu wahren sowie die weiteren sportrechtlichen Vorgaben nach den jeweils geltenden
Verbandsrichtlinien bei sportlichen Aktivitäten zu beachten.
\subsection{}
Die verschiedenen Arten der Mitgliedschaft bedingen die im Folgenden dargestellte Rechte und Pflichten:
\begin{itemize}
\item
Volljährige natürliche Personen:
    \begin{itemize}
        \item
        Alle Rechte und Pflichten.
    \end{itemize}
\item
Juristische Personen:
    \begin{itemize}
        \item
        Alle Rechte und Pflichten mit Ausnahme der Art der Teilnahmen im Probejahr, diese können
        beliebig sein, und der Art der verpflichtenden Haftpflichtversicherung, wenn diese aus der
        Natur der Person bedingt keine private Haftpflichtversicherung sein kann muss ein der Natur
        der juristischen Person entsprechendes Äquivalent bestehen.
    \end{itemize}
\item
Jugendliche (ab dem 16. Lebensjahr):
    \begin{itemize}
        \item
        Alle Rechte und Pflichten mit Ausnahme der Pflicht zur Entrichtung von Umlagen, jedoch
        eingeschränkt durch die gesetzlichen Einschränkungen welche sich aus der Minderjährigkeit
        der Mitglieder ergeben.
    \end{itemize}

\item 
Ehrenmitglieder:
        \begin{itemize}
            \item
            Alle Rechte und keine Pflichten mit Ausnahme der unter §3 Absatz 1 genannten Pflichten,
            welche nicht Mitgliedsbeiträge oder Umlagen betreffen.
        \end{itemize}
\item 
Fördermitglieder:
        \begin{itemize}
        \item
             Keine Rechte und keine Pflichten mit Ausnahme der Pflicht zur Zahlung des Förderbeitrages,
             der Förderung des Vereins, der Wahrung des Vereinsfriedens und des Rechtes auf Erhalt des
             jährlichen „Fördererbriefes“.
        \end{itemize}
        \item 
        Jugendliche und Kinder bis zu Vollendung des 16. Lebensjahres
        \begin{itemize}
            \item
        Keine Rechte und keine Pflichten. Die Zahlung des Kinderbeitrages und die weiteren Pflichten
        aus §3 Absatz 1 gehen in den Verantwortungsbereich der Eltern oder die gesetzliche
        Vormundschaft innehabende Person über.
        \end{itemize}
\end{itemize}
Die Regularien aus §3 Absatz 5 bzgl. des Austrittes und Ausschlusses von Mitgliedern bleiben hiervon unberührt
und gelten für jede Form der Mitgliedschaft.
\subsection{}
Die Aufnahme in den Verein ist schriftlich zu beantragen. Über den Aufnahmeantrag entscheidet die
Vorstandschaft. Die Ablehnung des Aufnahmeantrags kann dem Antragsteller ohne Angabe von Gründen
schriftlich per E-Mail oder auf dem Postweg mitgeteilt werden.
\subsection{}
Das 1. Jahr der Mitgliedschaft gilt, soweit unter §3 Absatz 2 nichts anderes angegeben, oder diese Pflicht
ausgenommen ist, als Probejahr. Innerhalb des Probejahrs sind das Mitglied und der Verein berechtigt, die
Mitgliedschaft mit Wirkung zum Ende des Monats, in dem die Kündigung ausgesprochen wurde, ohne Angaben
von Gründen zu kündigen. Die Probezeit kann durch den Verein einseitig durch die Vorstandschaft ohne Angabe
von Gründen einmalig um maximal 6 Monate verlängert werden.

In der einjährigen Probezeit muss das Mitglied an zwei Veranstaltungen teilnehmen, einmal als Teilnehmer und
einmal als Helfer. Während des Zeitraums der Mitgliedschaft hat eine private Haftpflichtversicherung zu
bestehen.
\subsection{Beendigung der Mitgliedschaft}
Die Mitgliedschaft endet mit dem Tod des Mitglieds, durch freiwilligen Austritt, Ausschluss aus dem Verein
oder Verlust der Rechtsfähigkeit der juristischen Person.
\subsubsection{Freiwilliger Austritt} Der freiwillige Austritt erfolgt durch eine schriftliche Erklärung gegenüber einem vertretungsberechtigten
Vorstandsmitgliedes. Er ist nur unter Einhaltung einer Kündigungsfrist von einem Monat zum Ende eines
Kalenderjahres zulässig. Die Mitgliedschaft ist nicht übertragbar.
\subsubsection{Ausschluss aus dem Verein}
Der Ausschluss aus dem Verein und die Streichung von der Mitgliedsliste erfolgt
\begin{itemize}
\item wenn das Mitglied trotz dreimaliger Mahnung an die zuletzt bekannte Adresse
länger als drei Monate mit seiner fälligen Beitragszahlung in Verzug ist, ohne 
dass eine soziale Notlage nachgewiesen wird. Bei einer sozialen Notlage kann die
Vorstandschaft die Beitragszahlungen ganz oder teilweise aufheben,
\item bei grobem Verstoß gegen die Satzung, Vereinsordnung oder
Verbandsrichtlinien,
\item wegen massiven unfairen oder unkameradschaftlichen Verhaltens,
\item wegen begründeter Bedenken zur Wahrung des Vereinsfriedens,
\item wegen unehrenhaften Verhaltens innerhalb und außerhalb des Vereinslebens,
wenn hierdurch die Interessen und das Ansehen des Vereins in der Öffentlichkeit
oder vereinsintern schwerwiegend beeinträchtigt werden.
\end{itemize}

Der Ausschluss eines Mitglieds aus dem Verein wird durch die Vorstandschaft mit einfacher Stimmenmehrheit
beschlossen und dem betreffenden Mitglied schriftlich per eingeschriebenen Brief mitgeteilt. Hiergegen kann das
Mitglied innerhalb eines Monats nach Zugang des Ausschlussschreibens schriftlich Berufung an die
Mitgliederversammlung einlegen. Diese entscheidet mit einfacher Mehrheit endgültig. Bis zum Abschluss dieses
vereinsinternen Verfahrens ruhen sämtliche Rechte des Mitglieds.
Mit Beendigung der Mitgliedschaft erlöschen alle Ansprüche aus dem Mitgliedsverhältnis, unbeschadet des
Anspruchs des Vereins auf bestehende Forderungen.
\section{Mitgliedsbeiträge / Umlagen}
Von den Mitgliedern werden, je nach Art der Mitgliedschaft, Beiträge erhoben, deren Höhe und Fälligkeit von
der Gesamtvorstandschaft mit einfacher Stimmenmehrheit beschlossen werden (Finanzordnung in der
Vereinsordnung). Von Vereinsmitgliedern, die Mitglied in mehreren Abteilungen sind, wird der Vereinsbeitrag
nur einmal erhoben.
Zur Abdeckung besonderer finanzieller Aufwendungen können nach Zustimmung durch die
Mitgliederversammlung Umlagen erhoben werden.
Die jeweils aktuelle Finanzordnung der Vereinsordnung kann weitere Beitragszahlungen / Umlagen für die
Zugehörigkeit zu einzelnen Abteilungen vorsehen.
Ehrenmitglieder sind grundsätzlich von der Beitrags- und / oder Umlagepflicht befreit.
\section{Rechte der Mitglieder}
Mitglieder, ausgenommen Fördermitglieder, ab der Anerkennung der staatlichen Volljährigkeit haben das aktive
Stimmrecht in der Mitgliedsversammlung und in den jeweiligen Abteilungsversammlungen.
Die stimmberechtigten Mitglieder haben das Recht, der Vorstandschaft und zur Mitgliedsversammlung Anträge
zu unterbreiten. Anträge zu Satzungsänderungen müssen der Vorstandschaft sechs Wochen vor der
Mitgliederversammlung eingereicht werden. Im Übrigen gilt § 7 dieser Satzung.
Alle Mitglieder werden bei der Teilnahme an Veranstaltungen des Vereins und Benutzung eventuell vorhandener
Übungsstätten unter Beachtung der Platz-, Hallen- beziehungsweise Hausordnung sowie sonstiger Ordnungen
bevorzugt.
Sie wählen die Vorstandsmitglieder und die jeweiligen Abteilungsleiter:innen, siehe §§ 10 und 13 dieser
Satzung. Eine Übertragung des Stimmrechts ist ausgeschlossen.
\section{Organe des Vereins}
Zu den Organen des Vereins gehören:
\begin{enumerate}
\item
Mitgliederversammlung
\item
Die Vorstandschaft
\end{enumerate}

\section{Mitgliederversammlung}
Die Mitgliederversammlung ist das oberste Organ des Vereins. Ihre Beschlüsse sind für alle Mitglieder und
Organe bindend. Die Mitgliederversammlung hat das Recht, gefasste Beschlüsse wieder aufzuheben.


Die ordentliche Mitgliederversammlung ist einmal im Jahr von der Vorstandschaft einzuberufen. Alle Mitglieder
sind unter Bekanntgabe der Tagesordnung und unter Einhaltung einer Frist von mindestens zwei Wochen
einzuladen.


Die Einladung hat schriftlich per E-Mail oder brieflich auf dem Postweg an die Mitglieder zu erfolgen.
Maßgeblich ist die letzte dem Verein mitgeteilte E-Mailadresse beziehungsweise Adresse.


Die Vorstandschaft kann jederzeit eine außerordentliche Mitgliederversammlung einberufen. Hierzu ist er
verpflichtet, wenn dies ein Viertel der stimmberechtigten Mitglieder schriftlich unter Angabe von Gründen
beantragt. In diesem Fall sind alle Mitglieder unter Bekanntgabe der Tagesordnung und unter Einhaltung einer
Frist von mindestens einer Woche einzuladen. Der obige Absatz 3 dieser Vorschrift (Einladungsvorgaben) gilt
entsprechend.

Anträge zur Mitgliederversammlung sind spätestens eine Woche vor der Mitgliederversammlung beim 1.
Vorsitzenden schriftlich per E-Mail oder auf dem Postweg einzureichen, soweit die Satzung keine andere Frist
vorschreibt.

Die Berücksichtigung verspäteter schriftlicher Anträge zu Mitgliederversammlungen ist nur möglich,
wenn jeder einzelne Antrag von mindestens der Hälfte der anwesenden stimmberechtigten Vereinsmitglieder durch
Beschluss nach Abstimmung zugelassen wird und der Antrag keine qualifizierte Mehrheit verlangt (Dringlichkeitsantrag).

Die Mitgliederversammlung ist beschlussfähig unabhängig von der Zahl der erschienenen Mitglieder.
\section{Aufgaben der Mitgliederversammlung}
Die Mitgliederversammlung ist für folgende Angelegenheiten zuständig:
\begin{itemize}
\item
die Wahl der Vorstandsmitglieder;
\item
die Entgegennahme des Jahres- und Kassenberichts der Vorstandschaft und des
Berichts der Kassenprüfer:innen sowie die Erteilung der Entlassung;
\item
die Wahl von bis zu zwei Kassenprüfer:innen;
\item
die Ernennung von Ehrenmitgliedern;
\item
die Beschlussfassung über Satzungsänderungen und alle sonstigen ihr von der
Vorstandschaft unterbreiteten Anträge;
\item
weitere Aufgaben, soweit sich dies aus der Satzung, Vereinsordnung oder dem
Gesetz ergibt.
\end{itemize}
\section{Beschlussfassung der Mitgliederversammlung}
Die Mitgliederversammlung fasst ihre Beschlüsse mit einfacher Mehrheit, soweit sich aus der Satzung
(insbesondere § 5) oder dem Gesetz nichts anderes ergibt. Eine Vertretung zur Stimmabgabe ist unzulässig.
Juristische Personen können sich durch ausgewiesene vertretungsberechtigte Personen vertreten lassen.

Satzungsänderungen und die Auflösung des Vereins bedürfen einer Mehrheit der Stimmen von drei Vierteln der
anwesenden stimmberechtigten Vereinsmitglieder.

Die Beschlussfassung erfolgt durch offene Abstimmung, auf Antrag eines der anwesenden Mitglieder erfolgt sie
in geheimer Abstimmung. Die Stimmenthaltung bleibt außer Betracht.

Kommt es bei der Wahl der Vorstandsmitglieder oder bei der Wahl der Kassenprüfer:innen zu Stimmgleichheit,
so findet eine Stichwahl statt. Bringt auch diese keine Mehrheit für eine Person, so wird durch Los entschieden.

Den Vorsitz in der Mitgliederversammlung führt der 1. Vorsitzende, bei dessen Verhinderung der 2.
Vorsitzende. Der Vorsitzende ist berechtigt, für einzelne Tagesordnungspunkte den Vorsitz / die Leitung an eine
andere Person zu übertragen.
\section{Die Vorstandschaft}
Die Vorstandschaft besteht aus dem 1. und dem 2. Vorsitzenden. Sie vertreten den Verein im Sinne des § 26
BGB gerichtlich und außergerichtlich und zeichnen als gesetzliche Vertreter:innen. Beide Vorstandsmitglieder
sind einzeln vertretungsberechtigt. Eine erweiterte Vorstandschaft ist für die Wahrung der Geschäftstüchtigkeit
des Vereins nicht zwingend erforderlich.
\section{Wahl der Vorstandschaft}
Die Vorstandschaft wird von der Mitgliederversammlung für die Dauer von einem Jahr gewählt. Sie bleibt bis
zur Neu- bzw. Wiederwahl der nachfolgenden Vorstandschaft im Amt. 

Scheidet ein Vorstandsmitglied nach § 26 BGB während der Amtsperiode vorzeitig aus, so bestimmt die
erweiterte Vorstandschaft durch Beschluss mit einfacher Mehrheit ein kommissarisches Vorstandsmitglied bis
zur nächsten Mitgliederversammlung.

Mit Beendigung der Mitgliedschaft im Verein endet automatisch auch das Amt als Vorstandsmitglied.
\section{Aufgaben und Zuständigkeit der Vorstandschaft}
Der Vorstandschaft sind alle Aufgaben des Vereins übertragen, die nicht satzungsgemäß in die Zuständigkeit
anderer Vereinsorgane fallen. Die Vorstandschaft kann intern eine Aufgaben- und Zuständigkeitsregelung
festlegen, die in der Vereinsordnung festzuhalten ist. Der Vorstandschaft obliegt insbesondere der Umgang mit
Behörden und Verbänden, die Entscheidung über alle Vertragsabschlüsse, deren Änderung und Kündigung
sowie alle weiteren rechtsgeschäftlichen Verpflichtungen.

Zur Zuständigkeit der Vorstandschaft gehören:
\begin{itemize}
\item
Verfassen und Beschluss der Vereinsordnung;
\item
Entscheidung über die Aufnahme neuer Mitglieder;
\item
Durchführung der Beschlüsse der Mitgliederversammlung;
\item
Delegation von Aufgaben und Einsetzen von Ausschüssen;
\item
Planung und Durchführung von Veranstaltungen;
\item
Vorprüfung der Gewinn- und Verlustrechnung, Haushaltsansätze, Finanzplanung;
\item
Schlichtung aller Streitigkeiten innerhalb des Vereins und Entscheidung über alle erhobenen
Widersprüche;
\item
Zusammenarbeit mit der erweiterten Vorstandschaft und angeschlossenen Abteilungen.
\end{itemize}
\section{Abteilungen des Vereins}
Innerhalb des Vereins werden für die unterschiedlichen Aktivitäten gesonderte Abteilungen eingerichtet. Die
Abteilungen sind rechtlich unselbstständige Untergliederungen des Vereins. Aus der Mitgliedschaft in einer
Abteilung ergeben sich keine über diese Satzung hinausgehenden Rechte und Pflichten, wenn nicht im
Folgenden etwas anderes bestimmt ist. Mitglied einer Abteilung kann nur werden, wer zugleich Mitglied des
Vereins ist.

Die Mitgliederversammlung entscheidet über die Gründung und Auflösung von Abteilungen. Die Entscheidung
ergeht mit einfacher Mehrheit. Bei der Auflösung einer Abteilung ist zuvor die Beschlussfassung durch die
Mitgliederversammlung einzuholen; der Wille der betroffenen Abteilung ist in der Wahlentscheidung der
Mitgliederversammlung des Vereins zu berücksichtigen.
Jede Abteilung nimmt ihre Angelegenheiten eigenverantwortlich war, soweit nicht diese Satzung dem
entgegensteht oder eine andere Abteilung bzw. der Verein hiervon betroffen ist. In diesen Fällen regelt die
Vorstandschaft unter Beachtung der einzelnen Belange die Angelegenheiten.

Die Leitung der Abteilung obliegt den jeweiligen Abteilungsleiter:innen und ggf. Stellvertreter:innen, die durch
die Mitglieder der Abteilung in einer einzuberufenden Abteilungsversammlung gewählt werden. Die Amtszeit
entspricht der satzungsgemäßen Amtszeit der Vorstandschaft. Sie bleiben bis zur Neu- bzw. Wiederwahl ihrer
Nachfolger:innen im Amt. Scheidet eine Person, welche die Abteilungsleitung innehat, vorzeitig aus und / oder
findet sich keine gewählte Person für die Position, so nimmt ein Mitglied der erweiterten Vorstandschaft die
Geschäfte der Abteilungsleitung zunächst kommissarisch war. Innerhalb eines Monats ist eine außerordentliche
Abteilungsversammlung einzuberufen, auf der eine neue Person zur Abteilungleitung durch die Mitglieder der
Abteilung für die noch ausbleibende Amtszeit zu wählen ist.

Die Leiter der Abteilungen sind besondere Vertreter:innen des Vereins im Sinne von § 30 BGB; sie können den
Verein beschränkt auf ihre Abteilung und beschränkt auf das Aktivvermögen der Abteilung bis zu einer Höhe
von 100 € rechtsgeschäftlich vertreten. Die Eingehung von Anstellungs-, Miet- oder Leasingverträgen oder
sonstige Verträge als Dauerschuldverhältnisse, bedürfen grundsätzlich der vorherigen Zustimmung der
Vorstandschaft.

Die Abteilungsleiter:innen haben, der Vorstandschaft in jeder Vorstandsitzung und im Bedarfsfall auch
außerhalb davon über Aktivitäten und Vorkommnisse in den Abteilungen zu unterrichten.

Die Abteilungen können sich eigene Abteilungsordnungen geben, die erst mit der Genehmigung der
Vorstandschaft wirksam werden. Die Abteilungsordnungen müssen die Organisation der Abteilung regeln und 
sich an den Vorgaben dieser Satzung orientieren.
Im Kollisionsfall hat die Vereinssatzung Vorrang, die weiterhin verbindlich für alle Mitglieder gilt. Über neue
oder geänderte Abteilungsordnungen ist die Mitgliederversammlung zu informieren.

Ein Vereinsmitglied kann Mitglied mehrerer Abteilungen sein. Über die Aufnahme in eine Abteilung entscheidet
die jeweilig die Abteilungsleitung innehabende Person. Ein Vereinsmitglied hat das Recht, jederzeit zwischen
den Abteilungen zu wechseln, soweit nicht bestehende Kapazitätsgrenzen dem entgegenstehen. Für diesen Fall
sind Wartelisten einzurichten. Die Kapazitätsgrenzen werden durch die erweiterte Vorstandschaft nach
Anhörung der Abteilungsleitung festgelegt.

Der Vereinsführung obliegt ansonsten die Mitgliederverwaltung. Soweit für die Organisation erforderlich, kann
jede Abteilung von der zentralen Mitgliederverwaltung Listen über ihre Abteilung erhalten.

Soweit erforderlich, erwirbt der Verein die Mitgliedschaft in Fachverbänden; die daraus resultierenden Rechte
und Pflichten erstrecken sich auch auf die Mitglieder der Abteilungen.
s
Die Regelungen der §§ 5, 7 bis 9 dieser Satzung gelten, soweit möglich, entsprechend für die
Abteilungsversammlungen.
\section{Kassenprüfer:in}
Der/Die Kassenprüfer:in wird von der Mitgliederversammlung für die Dauer von einem Jahr gewählt. Diese
Person darf nicht Mitglieder der erweiterten Vorstandschaft sein. Sie hat das Recht, die Vereinskasse und die
Buchführung jederzeit zu überprüfen. Über die Prüfung der gesamten Buch- und Kassenführung hat sie der
Mitgliederversammlung Bericht zu erstatten. Das Prüfungsrecht erstreckt sich nur auf die buchhalterische
Richtigkeit, nicht auf die Zweckmäßigkeit der Vorgänge. Sie kann nur einmal wiedergewählt werden.
\section{Protokollierung}
Der Verlauf der Mitgliederversammlungen, der Vorstandssitzungen sowie die Abteilungstreffen sind zu
protokollieren. Das Protokoll der Mitgliederversammlung wird von der schriftführenden Person und der
Versammlungsleitung unterzeichnet. Die Protokolle der Vorstandssitzungen sind von einem
vertretungsberechtigten Mitglied der Vorstandschaft abzuzeichnen. Die Vorstandsprotokolle, die Protokolle der
Abteilungstreffen und die der Mitgliederversammlung hat die Vorstandschaft aufzubewahren.
Die einzelnen Abteilungen sind selbstständig dafür zuständig, dass ein Protokoll bei einem Abteilungstreffen
erstellt und zur Vorstandschaft geschickt wird.
\section{Auflösung des Vereins}
Über die Auflösung des Vereins beschließt eine zu diesem Zwecke besonders einberufene
Mitgliederversammlung mit einer Mehrheit von drei Vierteln der anwesenden stimmberechtigten Mitglieder,
vorausgesetzt mindestens ein Viertel aller stimmberechtigten Vereinsmitglieder ist anwesend.
Ist diese Zahl nicht erreicht, muss innerhalb von zwei Wochen eine neue Mitgliederversammlung von der
Vorstandschaft einberufen werden, die alsdann mit einer Mehrheit von drei Vierteln der anwesenden
stimmberechtigten Vereinsmitglieder die Auflösung beschließt.

Die Mitgliederversammlung ernennt zur Abwicklung der Geschäfte Liquidatoren.

Bei Auflösung des Vereins oder bei Wegfall steuerbegünstigter Zwecke fällt das Vermögen des Vereins an eine
juristische Person des öffentlichen Rechts oder eine andere steuerbegünstigte Körperschaft zwecks Verwendung
für die Förderung von Kunst und Kultur, insbesondere der Erhalt und die Pflege des historischen Brauchtums des
Mittelalters, der mittelalterlichen Musik und des Tanzens, des mittelalterlichen Schwertfechtens und auch für
weitere Betätigungsfelder des lebendigen Mittelalters.

Wird mit der Auflösung des Vereins nur eine Änderung der Rechtsform oder eine Verschmelzung mit einem
anderen gleichartigen Verein angestrebt, wobei die unmittelbare ausschließliche Verfolgung des bisherigen
Vereinszweckes durch den neuen Rechtsträger weiterhin gewährleistet wird, geht das Vermögen an den neuen
Rechtsträger über.

Vor Durchführung der Auflösung und Weitergabe des noch vorhandenen Vermögens ist zunächst das Finanzamt
zu hören.
\end{document}
